\section{Introduction}
\label{sec:intro}

\IEEEPARstart{M}{obile} applications run neural networks on images, speech, and sensor data, but a full model can exceed the compute and memory budget of a phone. Split inference divides the model between the device and a nearby edge server. The device runs the first layers and sends an intermediate feature, and the edge runs the rest of the model~\cite{kang2017neurosurgeon,laskaridis2020spinn}. Split learning trains such a model across many users without collecting their raw data~\cite{gupta2018distributed,vepakomma2018split,thapa2022splitfed}. Personalized split learning adds a classifier at the end of the device part~\cite{han2023splitgp,rizwan2026splitomc}. We call each input for which an application needs a class label, such as one camera frame, a \emph{request}. Each request therefore has two classifiers available. The \emph{device exit} is personalized to the classes in its user's training data, and the \emph{edge exit} is trained across all users of a cell.

Which of the two classifiers gives the right answer depends on where the user is. Consider a camera application that recognizes objects. At home it mostly sees objects that the user has photographed before, and in a shopping district it sees products and signs that other users of that area have photographed. In the commute scenario that we simulate, about one in nine requests of a device at home comes from classes outside its own training data, and about one in two after the user leaves home. The two exits are accurate on different parts of this traffic. Over the first simulated day, the device exit is correct on 86\% of requests from the user's own classes but on only 7\% of requests from classes learned by other users of the current cell. The edge exit is correct on 72\% and 41\% of the same requests. Neither exit is the better choice for the whole day.

A natural answer is to choose an exit for each request. Common split inference answers on the device when the device exit is confident and otherwise sends the request to the edge~\cite{han2023splitgp,laskaridis2020spinn,teerapittayanon2016branchynet}. This rule works well when the device exit is uncertain on classes it has not learned. A personalized classifier often is not. It places most of its probability on the user's own classes even for a request from another class, so its entropy separates the two kinds of requests only weakly (an area under the ROC curve of 0.64). Confidence-based offloading therefore sends 92\% of the requests to the edge and is about as accurate as always using the edge exit.

Another answer is to change how strongly each device personalizes during training, for example by keeping more of the local model at home and more of the shared model while away~\cite{deng2020apfl}. We tested such a policy with an oracle that knows whether each device is at home. It raised the accuracy of devices at home by 3.7 points and lowered that of devices away from home by 1.7 points. Because the policy changes the personalization of residents and visitors together, this experiment shows a trade-off of the policy but does not isolate its cause. These observations led us to look at inference rules that leave training unchanged.

A device does not need to choose between its two exits. If it averages the class probabilities of both exits, the edge exit's answer remains large on a request that the device exit has never learned, as long as the edge exit is confident and the device exit spreads its probability over the user's classes. A product of probabilities, which is what adding logits computes, instead scales the edge probability by the device probability, so a near-zero device probability can outweigh a confident edge answer. The average still has a problem. The edge exit is trained on the class mix of the whole cell, in which the device's own classes make up about one quarter of the samples, while a device at home receives these classes in about 88\% of its requests. The edge exit therefore discounts the classes that dominate the device's traffic at home.

We present DriftGate, an inference method for personalized split learning under user mobility. DriftGate answers with a weighted average of the class probabilities of the two exits. Before averaging, it corrects the edge exit for the class mix of its cell, so that the device's own classes are no longer discounted. It sets the weight of each exit from how confident the two exits have been on the device's recent offloaded requests, without labels. Finally, it sends only the requests on which the device exit is uncertain to the edge, so the device can trade accuracy for communication. DriftGate does not change training and adds a few microseconds of computation per request.

We evaluate DriftGate in eight settings. They include a commute scenario, real mobility traces from the GeoLife dataset~\cite{zheng2010geolife}, faster movement, partial participation, cell-wide changes in the class mix, random mobility, 100 classes, and a ResNet-18 split~\cite{he2016resnet}. We compare it with ten inference rules that cover confidence-based offloading, each single exit, the probability average, the logit sum used by FedRoD~\cite{chen2022fedrod}, the entropy objective of FedTHE~\cite{jiang2023fedthe}, the geometric ensembling with early exit of ZTW~\cite{wolczyk2021ztw}, label-shift correction~\cite{saerens2002adjusting}, and a learned per-request weight. Because our scenarios train the models from scratch during one simulated day, we report three evaluation conditions: the whole first day, the rounds after the first hours of training, and a replay of the day with the trained models frozen. Over the first day, DriftGate has the highest mean accuracy in all eight settings, 0.1 to 1.9 points above the strongest alternative and 5.4 to 12.5 points above confidence-based offloading. After round 30, it ranks first in six settings, and the margins over the strongest alternative range from $-0.2$ to $+1.1$ points. In the replay with trained models, it is 0.6 points above the strongest alternative on GeoLife traces and 0.3 points below the probability average in the commute scenario. \todo{Summary of measured on-device latency and energy.}

This paper makes the following contributions.
\begin{itemize}
    \item We show how user mobility affects inference in personalized split learning. The device exit and the edge exit are accurate on different requests, and the confidence of the device exit separates these requests only weakly.
    \item We design DriftGate, which averages the probabilities of the two exits after correcting the class bias of the edge exit, sets the weight of each exit without labels, and supports selective offloading. It runs on top of the trained model and does not change training.
    \item We evaluate DriftGate against ten inference rules in eight settings, including real mobility traces, separate the effect of early training from that of trained models, and report which parts of the design matter under each condition.
\end{itemize}

The rest of the paper is organized as follows. Section~\ref{sec:motivation} measures how the two exits and the two common remedies behave under mobility. Section~\ref{sec:related} reviews related work. Section~\ref{sec:system} defines the system and the inference problem, and Section~\ref{sec:design} presents DriftGate. Section~\ref{sec:evaluation} reports the evaluation, Section~\ref{sec:discussion} discusses the scope of the method, and Section~\ref{sec:conclusion} concludes.
